% Part: lambda-calculus
% Chapter: introduction
% Section: overview

\documentclass[../../../include/open-logic-section]{subfiles}

\begin{document}

\olfileid{lam}{int}{ovr}
\olsection{అవలోకనం}

లాంబ్డా కలనశాస్త్రాన్ని అలోంజో చర్చ్ 1930ల తొలినాళ్లలో నిర్మాణాత్మక
తర్కానికి పునాదిగా మొదట రూపొందించారు; గణనీయ ప్రమేయాలకు ఒక నమూనాగా
మాత్రం \emph{రూపొందించలేదు}. కానీ అది ట్యూరింగ్-గణనీయ ప్రమేయాలు,
పాక్షిక పునరావృత్త ప్రమేయాలు వంటి ఇతర గణనీయతా నిర్వచనాలకు తుల్యమని
త్వరలోనే చూపించారు. మొదట ఇది కొంత ఆశ్చర్యాన్ని కలిగించిందనే వాస్తవం
ఈ లక్షణ నిర్ధారణను మరింత ఆసక్తికరంగా చేస్తుంది.

ఒక ప్రమేయానికి పేరు పెట్టి నిర్వచించడానికి బదులు, దాన్ని నిర్వచించే
సాంకేతిక వ్యక్తీకరణతో నేరుగా సూచించడానికి లాంబ్డా సంకేతనం అనుకూలమైన
మార్గం. ``$f(x) = x + 3$ ద్వారా నిర్వచించిన ప్రమేయం $f$ అనుకుందాం''
అని చెప్పడానికి బదులు, ``$f$ అనేది $\lambd[x][(x + 3)]$ ప్రమేయం
అనుకుందాం'' అని చెప్పవచ్చు. మరో మాటలో, తన ఆర్గ్యుమెంట్‌కు మూడును
కలిపే ప్రమేయానికి $\lambd[x][(x+3)]$ కేవలం ఒక \emph{పేరు}.
ఈ వ్యక్తీకరణలో $x$ ఒక డమ్మీ చరం, అంటే స్థానపూరకం: అదే ప్రమేయాన్ని
$\lambd[y][(y + 3)]$ అని కూడా సూచించవచ్చు. ఇతర పరామితులు ఉన్నప్పటికీ
ఈ సంకేతనం పనిచేస్తుంది. ఉదాహరణకు, $g(x, y)$ రెండు చరాల ప్రమేయమనీ,
$k$ ఒక సహజ సంఖ్య అనీ అనుకుందాం. అప్పుడు ఏ~$x$నైనా~$g(x, k)$కు
పంపే ప్రమేయమే $\lambd[x][g(x,k)]$.

సాంకేతిక వ్యక్తీకరణ నుంచి ప్రమేయాన్ని ఇలా నిర్వచించడాన్ని
\emph{లాంబ్డా అమూర్తీకరణ} అంటారు. లాంబ్డా అమూర్తీకరణకు మరో వైపు
\emph{ప్రయోగం}: ఒక ప్రమేయం $f$ (ఉదాహరణకు, సహజ సంఖ్యలపై నిర్వచించినది)
ఉంటే, దాన్ని 2 వంటి ఏ విలువకైనా \emph{ప్రయోగించవచ్చు}. సాంప్రదాయ
సంకేతనంలో ఫలితాన్ని~$f(2)$ అని రాస్తాం.

లాంబ్డా అమూర్తీకరణను ప్రయోగంతో కలిపితే ఏమవుతుంది? అమూర్తీకరించిన
చరం స్థానంలో ప్రయోగించబడే పదాన్ని ``ప్రవేశపెట్టడం'' ద్వారా ఫలిత
వ్యక్తీకరణను సరళీకరించవచ్చు. ఉదాహరణకు,
\[
(\lambd[x][(x + 3)])(2)
\]
ను~$2 + 3$కు సరళీకరించవచ్చు.

ఇప్పటివరకు మనం సాంప్రదాయ భావనలకు కొత్త సంకేతనాలను ప్రవేశపెట్టడం తప్ప
మరేమీ చేయలేదు. అయితే లాంబ్డా కలనశాస్త్రం సమితి-సిద్ధాంత దృక్కోణం నుంచి
మరింత మౌలికమైన మలుపును సూచిస్తుంది. ఈ చట్రంలో:
\begin{enumerate}
\item ప్రతిదీ ఒక ప్రమేయాన్ని సూచిస్తుంది.
\item లాంబ్డా అమూర్తీకరణతో ప్రమేయాలను నిర్వచించవచ్చు.
\item దేనినైనా మరే దానికైనా ప్రయోగించవచ్చు.
\end{enumerate}
ఉదాహరణకు, $F$ లాంబ్డా కలనశాస్త్రంలోని ఒక పదమైతే, $F(F)$కు ఎల్లప్పుడూ
అర్థం ఉందని భావిస్తాం. ఈ ఉదారమైన చట్రాన్ని \emph{టైపులేని} లాంబ్డా
కలనశాస్త్రం అంటారు; ఇక్కడ ``టైపులేని'' అంటే ``దేనిని దేనికి
ప్రయోగించవచ్చనే దానిపై ఎలాంటి ఆంక్ష లేదు'' అని అర్థం.

\begin{digress}
టైపులేని రూపానికి ముఖ్యమైన భేదరూపమైన \emph{టైపుగల} లాంబ్డా
కలనశాస్త్రం కూడా ఉంది. అనేక విధాలుగా టైపుగల లాంబ్డా కలనశాస్త్రం
టైపులేని రూపాన్ని పోలి ఉన్నప్పటికీ, దాన్ని సాంప్రదాయ సమితి-సిద్ధాంత
చట్రంతో సమన్వయించడం చాలా సులభం; దానికి చాలా భిన్నమైన కొన్ని లక్షణాలు
కూడా ఉన్నాయి.

లాంబ్డా కలనశాస్త్రంపై పరిశోధన సైద్ధాంతిక కంప్యూటర్ శాస్త్రంలోనూ,
ప్రోగ్రామింగ్ భాషల రూపకల్పనలోనూ కేంద్రస్థానంలో నిలిచింది. 1950లలో
జాన్ మెకార్తీ రూపొందించిన LISP ఈ భావనల ప్రభావానికి లోనైన భాషకు ఒక
తొలి ఉదాహరణ.
\end{digress}

\end{document}
